\documentclass[11pt,a4paper]{article}
\usepackage[margin=27mm]{geometry}
\usepackage{amsmath,amssymb,amsthm,booktabs,hyperref}
\hypersetup{hidelinks}
\newtheorem{lemma}{Lemma}
\newcommand{\Aut}{\operatorname{Aut}}
\newcommand{\Soc}{\operatorname{Soc}}
\title{Ternary modules on the flags of the plane of order four}
\author{Private derivation and exact certificate: independent review pending}
\date{}
\begin{document}\maketitle

This bounded note treats L050, Theorem~1(a), Proposition~1,
Corollary~1 and Remarks~1--2. The source is Leemans--Rodrigues,
\emph{Linear codes with complementary duals from some strongly regular
subgraphs of the McLaughlin graph}, Math. Commun.\ 21 (2016), 239--249.
We reuse the common matrix-algebra, minimum-support and transitive
information-set arguments. Complete previously accepted enumerations
are retained for the two main codes. The two additional codes in
Remark~2 require only the bounded search proved below.

\section{The graph and the two principal codes}
Use $\mathbb F_4=\mathbb F_2[z]/(z^2+z+1)$ and the $105$ incident
point--line flags of $\operatorname{PG}(2,4)$. Join distinct flags
$(p,\ell),(q,m)$ when $p=q$, or $\ell=m$, or neither cross incidence
$p\in m$, $q\in\ell$ holds. The complete field construction checks
parameters $(105,72,51,45)$.

The printed adjacency sentence on page 242 does not state this condition
correctly. Reading its clauses conjunctively gives valency eight;
reading them all disjunctively gives valency forty. Our identification
with the intended graph therefore retains the uniqueness of the strongly
regular graph with parameters $(105,72,51,45)$, obtained from uniqueness
of its complement in Coolsaet's Theorem~3, cited as source reference [6].
The parameter check alone would not identify an arbitrary graph with a
named source graph without this input.

Let $A$ be the adjacency matrix over $F=\mathbb F_3$, and write
\[
 C=\langle\text{rows of }A\rangle,
 \qquad D=\langle A_i-A_1:1\leq i\leq105\rangle.
\]
The parameter identity gives $A^2=0$, so $C$ and $D$ are self-orthogonal.
Exact row reduction gives $\dim C=20$, $\dim D=19$, and
\[
 C=D\oplus Fj,                                      \tag{1}
\]
where $j$ is the constant vector. Both bases agree byte for byte with the
separately completed enumeration. The two complete distributions give
\[
\begin{array}{c|rrrr}
 &\dim&d&\text{minimum words}&(\dim C^\perp,d^\perp,\text{minimum words})\\\hline
 C&20&33&770&(85,6,15680)\\
 D&19&36&3360&(86,5,84).
\end{array}
\]
The complete dual distributions are checked again by the classical
MacWilliams identity. No $3^{20}$ or $3^{19}$ enumeration is repeated.
The retained distributions agree entry by entry with the two printed
tables.

\section{The full group and irreducibility}
The literal coordinate permutations in the certificate preserve every
entry of $A$. Their closure has order $241920$, and their action has
rank four. They generate a group $G$ with a normal subgroup
$N\cong\operatorname{PSL}_3(4)$ of order $20160$ and an explicit
complement $Q\cong D_{12}$ of order twelve. The complement consists of
a diagonal projectivity, field conjugation and point--line duality.
The checker derives their actual flag permutations, verifies
$N\cap Q=1$ and $NQ=G$, and checks the group identifications.
Thus $G\cong L_3(4):D_{12}$, with this action fixed.

The action on the row-reduced $19$-dimensional basis of $D$ is checked
directly. The generator matrices span a matrix algebra of dimension
$19^2=361$: the verifier evaluates independent generator words and
row-reduces their flattened matrices. A nonzero invariant subspace stable
under all these matrices is stable under every endomorphism, hence equals
$D$. Independence persists over field extensions, so $D$ is absolutely
irreducible. Consequently (1) is an internal decomposition into two simple
modules, as required in Proposition~1.

The complete minimum-word set of $C$ is obtained as four disjoint orbits
of literal weight-$33$ words, of lengths $105,105,280,280$.
The accepted count $770$ proves exhaustion, and also recovers Remark~1.
There are $385$ distinct supports. As for the $120$- and $176$-coordinate
codes, form their bipartite coordinate--support incidence graph.
Refine vertex colours by old colour and neighbour-colour multiplicities,
keeping the two parts initially distinct. Successively individualizing
coordinates $1,2,3,4,7,12,6$ gives cell sizes
\[
 105,8,3,2,4,3,4.
\]
All coordinate colours then differ. At each step a stabilizer orbit is
contained in the selected colour class; after the final step the
stabilizer is trivial. Thus
\[
 |\Aut(C)|\leq105\cdot8\cdot3\cdot2\cdot4\cdot3\cdot4=241920.
\]
Since $G\leq\Aut(\Gamma)\leq\Aut(C)$, equality follows.
Every automorphism of $D$ fixes $j$ and hence preserves $D+Fj=C$.
Conversely $G$ preserves $D$, so
\[
 \Aut(D)=\Aut(C)=\Aut(\Gamma)=L_3(4):D_{12}.          \tag{2}
\]
These are coordinate-permutation groups, not monomial groups.

Both retained distributions have exactly $3360$ words of weight $36$.
Since $D\subset C$, all such words of $C$ are in $D$.
Their span is a nonzero $G$-submodule of the irreducible $D$, and is
therefore $D$. This establishes the source's intrinsic definition of its
codimension-one code, rather than identifying it only by parameters.

\section{Uniqueness inside the permutation module}
Put $M=F^{105}$. A supplied basis gives invariant prefixes of dimensions
\[
 0,1,20,39,84,85,104,105.
\]
The factor dimensions are $1,19,19,45,1,19,1$.
All factors are simple. For the factors of dimensions one and nineteen,
this is verified by the full matrix algebra as above. To avoid an
unnecessarily large elimination for the $45$-dimensional factor, the
certificate uses the following equivalent sufficient test.

\begin{lemma}\label{lem:rankone}
Let $\mathcal A$ be an algebra generated by matrices on $F^d$.
Suppose it contains a rank-one matrix $P=vu$, and that the vectors $gv$
and $uh$, with $g,h$ words in its generators, each span their respective
column and row spaces of dimension $d$. Then $\mathcal A=M_d(F)$.
\end{lemma}
\begin{proof}
Choose $d$ such columns and $d$ such rows forming bases. Their $d^2$
outer products are a basis of $M_d(F)$; each is $gPh\in\mathcal A$.
\end{proof}

Here a recorded polynomial of a recorded sum of generator words gives
$P$ of rank one. The checker evaluates it and produces both bases of
size $45$. Hence its algebra has dimension $2025$, proving simplicity
over $F$ and after extension. This is a certificate of the concrete
factor, not an assumed classification of representations.

For each of the three $19$-dimensional factors $S_i$, all linear equations
\[
 XP_g=(S_i)_gX
\]
are solved, where $P_g$ is the actual $105$-coordinate permutation
matrix. The full solution-space dimensions are $0,1,1$ in the displayed
filtration order. Every basis solution has all rows in $D$, hence every
homomorphism $S_i\to M$ has image in $D$.

Any $19$-dimensional submodule of $M$ is simple. Indeed its composition
factors occur among those of $M$; dimension $45$ is impossible, while
the three available one-dimensional factors alone cannot account for
dimension nineteen. Its sole factor must therefore have dimension
nineteen. Such a simple submodule is isomorphic to one of the $S_i$,
and its inclusion is the image of a corresponding homomorphism into
$M$. The preceding complete equations force its image to be $D$.
Thus $D$ is the unique submodule of dimension nineteen in this
permutation representation, proving the first assertion of Remark~2.

\section{The other two twenty-dimensional modules}
The certificate supplies invariant row-reduced bases for $C_1,C_2$,
besides $C_3=C$. Their containments are
\[
\begin{array}{c|cc}
 &D\subseteq C_i&j\in C_i\\\hline
C_1&\text{no}&\text{yes}\\
C_2&\text{yes}&\text{no}\\
C_3&\text{yes}&\text{yes}.
\end{array}
\]
The only one-dimensional submodule of $M$ is $Fj$. To see this directly,
an invariant line has a spanning vector with entries in $\{0,a,-a\}$.
Transitivity makes its zero set empty. A nontrivial scalar action would
interchange the two nonzero values, forcing their multiplicities to be
equal, contrary to the odd length $105$. Thus the line is fixed and
constant. Every $20$-dimensional submodule has factors of dimensions
$19$ and one, by the full composition series above. The containments
and uniqueness just proved give
\[
 \dim\Soc(C_1)=1,\quad \dim\Soc(C_2)=19,\quad
 \dim\Soc(C_3)=20.
\]
The three modules are therefore reducible and pairwise nonisomorphic
as $FG$-modules, not just inequivalent as codes.

For their remaining code parameters, we reuse the already proved
transitive information-set argument. If a transitive group preserves a
length-$n$, dimension-$k$ code and $I$ is an information set, averaging
over the group shows that a word of weight $w$ has a translate whose
restriction to $I$ has weight at most $\lfloor kw/n\rfloor$.
For $(n,k)=(105,20)$ and candidate distances $30,28$, both the smaller-word
and minimum-word tests require only messages of weight at most five.
There are
\[
 \sum_{i=0}^{5}\binom{20}{i}2^i=583569
\]
such ternary messages, including zero. The complete shared enumerator
checks them once each. Its smallest nonzero output weights are $30,28$,
and literal witnesses attain those values. The averaging theorem
therefore proves those minimum distances.

The unions of a literal witness orbit and its negative have sizes $672$
and $720$. They are invariant sets of minimum words. Exactly $342$ and
$396$ of these words, respectively, have message weight at most five;
these equal the complete restricted enumeration counts at the minimum
weight. Thus every restricted minimum word is in the proposed invariant
set. Averaging an arbitrary minimum word into that restriction proves
that the sets are complete. This recovers
$C_1=[105,20,30]_3$ with $672$ minimum words and
$C_2=[105,20,28]_3$ with $720$ minimum words, as printed in Remark~2.
It neither assumes a complete $3^{20}$ enumeration nor asserts that a
short search is the full weight distribution.

\section{Qualifications and attribution}
The source's proof of Corollary~1 does not justify its global statement
that nineteen is the smallest dimension of a nontrivial irreducible
$F_3G$-module. Literally that statement is false: $G$ has a normal
subgroup of index two, so its quotient acts nontrivially by the scalars
$1,-1$ on a one-dimensional $F_3$-space. The certificate exhibits the
index-two subgroup by adjoining the diagonal and field automorphisms
to $N$; its order is $120960$. No faithfulness hypothesis is stated in
the printed claim. We prove the precise local replacement: the unique
nonconstant simple submodule of dimension nineteen in this permutation
module is $D$, and no nonconstant simple submodule has smaller dimension.
No classification of all faithful irreducible representations is inferred.

The common refinement and matrix-algebra arguments are the same as in
the earlier two code supplements. The averaging theorem and complete
sparse-message enumerator are reused from the distinctly reviewed
transitive-distance packet. The exact incidence geometry is classical,
and the uniqueness input identifying the source graph remains explicitly
retained. Source comparisons to historical optimal-code tables remain
historical assertions. Discovery used GRAPE and MeatAxe; the final
certificate checks literal generators, a full invariant filtration,
matrix products and complete linear equations, without requiring their
discovery outputs to be trusted. No whole-source admission, new source
credit, historical priority, numerical compression or end-to-end Lean
claim is made in this bounded packet.
\end{document}
